\chapter{Fuentes y Modelo de Datos}
\label{cap:datos}

\section{El Horario Programado}

El horario oficial se publica en GTFS y se transforma, una vez al día, en una base de datos SQLite de solo lectura
por red. La Tabla~\ref{tab:esquema} resume las tablas que usa el motor; el esquema completo está en el
Apéndice~\ref{ap:formatos}.

\begin{table}[ht]
\centering
\caption{Tablas de la Base de Datos del Horario que Usa el Motor}\label{tab:esquema}
\begin{tabularx}{\textwidth}{@{}lX@{}}
\toprule
Tabla & Contenido relevante para el motor \\
\midrule
\texttt{stops} & Paradas: identificador, nombre, latitud y longitud. \\
\texttt{service\_journeys} & Viajes: línea, recorrido (\texttt{journey\_pattern\_id}), número de viaje, calendario,
  hora de salida de cabecera y trazado (\texttt{shape\_id}). \\
\texttt{passing\_times} & Para cada viaje y cada parada en orden (\texttt{seq\_order}, desde 1): llegada y salida
  programadas en segundos y \texttt{dist\_m}, la distancia a lo largo del trazado. \\
\texttt{shapes} & Trazado de cada recorrido como lista JSON de puntos $[\varphi,\lambda,\sigma]$: latitud, longitud y
  distancia recorrida. \\
\texttt{trip\_aliases} & Correspondencia entre cada \texttt{trip\_id} original del GTFS y el viaje fusionado al que
  pertenece. \\
\texttt{journey\_patterns} & Recorridos: secuencias de paradas compartidas por varios viajes. \\
\bottomrule
\end{tabularx}
\end{table}

La columna \texttt{dist\_m} de \texttt{passing\_times} es el campo
\texttt{shape\_dist\_traveled} de \texttt{stop\_times.txt}, y la tercera componente de cada punto del trazado es el
\texttt{shape\_dist\_traveled} de \texttt{shapes.txt}. Ambas se miden sobre el mismo trazado y en la misma unidad (o
se corrigen al construir la base hasta que lo hacen; sección «\nameref{sec:trazados-gtfs}»), lo que permite pasar de un punto del trazado a la parada que tiene delante o detrás (sección «\nameref{sec:trazado}»). Los
puntos del trazado se guardan con 7 decimales en las coordenadas y 2 en la distancia. Si la red no tiene trazados, la
base no tiene la tabla \texttt{shapes} y el motor trata \texttt{dist\_m} como desconocida; esto tiene consecuencias en
el respaldo por longitud (sección «\nameref{sec:respaldo-largo}»).

\section{Trazados que No Casan con sus Paradas}
\label{sec:trazados-gtfs}

Lo anterior supone que el trazado y las distancias de las paradas casan, y en el GTFS de Bizkaibus no siempre es así.
Con el horario del 3 de noviembre de 2026, 55 de los 440 trazados que usan los viajes están dibujados del final al
principio: la primera parada del viaje queda junto al último punto del trazado y la última junto al primero. Afectan a
670 de los 6\,130 viajes fusionados, repartidos en 25 líneas. En otros trazados el sentido es el bueno, pero los metros
de \texttt{shape\_dist\_traveled} no caen sobre la parada: unos van desplazados 200 o 300~m, y en otros la diferencia
crece hasta varios kilómetros al final del viaje. En total, 1\,417 viajes tenían alguna parada a más de 200~m del punto
del trazado que marcan sus metros.

El efecto en el motor es grave. La ventana de búsqueda (sección «\nameref{sec:ventana}») se define con los metros de
las paradas, así que, con el trazado al revés, el autobús se busca en la otra punta del recorrido: un autobús a 9~m de su
parada queda «a kilómetros» y el pasillo de 50~m lo descarta. Esos autobuses no se situaban con el GPS y, como el
aprendizaje sale de las posiciones, tampoco aprendían: los viajes con el trazado al revés tenían el 47\,\% de sus
tramos con tiempo aprendido, frente al 95\,\% del resto.

La construcción de la base corrige las dos cosas antes de publicarla (el programa, \texttt{scripts/build/shapes.php},
es parte del código privado de la aplicación; el algoritmo completo es este):

\begin{enumerate}
  \item \emph{Sentido}. Un trazado está al revés si, en la mayoría de sus viajes, la suma de las distancias de la
        primera parada al último punto y de la última al primero es menor que la de cada parada a su extremo. Se le da
        la vuelta: los puntos en orden inverso y $\sigma'=\sigma_{\max}-\sigma$.
  \item \emph{Metros de las paradas}. Si el trazado se ha dado la vuelta, o si alguna parada cae a más de 100~m del
        punto del trazado que marcan sus metros, se miden de nuevo las de ese viaje: cada parada se proyecta sobre el
        trazado a partir de la anterior y, como mucho, hasta tres veces la distancia en línea recta entre las dos más
        500~m. Entre dos pasos del autobús por la misma zona gana el primero, salvo que el otro quede 30~m más cerca.
\end{enumerate}

\noindent Con el horario de ese día se dieron la vuelta los 55 trazados y se midieron de nuevo las paradas de 2\,905
viajes; los viajes con alguna parada a más de 200~m de sus metros bajaron de 1\,417 a 26. La corrección es general:
se aplica al construir cualquier red, porque nada en el formato GTFS impide estos errores.

Tampoco obliga a publicar los metros. El GTFS de Lurraldebus, Bilbobus, Tuvisa y La Unión no trae los de las paradas,
y el de Lurraldebus, La Unión y Bilbobus tampoco los del trazado (el de Bilbobus los deja a 0 en todos los puntos menos
el último). Sin metros, el motor sitúa al autobús con la recta entre paradas (sección «\nameref{sec:sin-trazado}»).
Por eso la construcción los mide: los del trazado, sumando la distancia de cada punto al siguiente, y los de las
paradas, como en el paso~2.

Los feeds de Lurraldebus y Bilbobus, las dos redes con motor, no traían vehículos el 3 de noviembre de 2026, así que
la medida se hizo con Bizkaibus: se le quitaron los metros del GTFS y se compararon tres versiones del mismo horario
con las mismas capturas, 36 bloques de una hora del 2 y el 3 de noviembre con 2,08 millones de llegadas de paso
conocido (tabla~\ref{tab:sin-metros}). Los metros medidos caen a una mediana de 21~m de los del GTFS, el 95\,\% a
menos de 103~m, y aciertan casi como ellos. La recta entre paradas sitúa con el GPS un 5,6\,\% menos de llegadas y
acierta 1,8 puntos menos dentro de $\pm$2~min. Sin tiempos aprendidos, como en una red recién activada, las tres
versiones quedan a menos de medio punto: entonces lo que falta es aprender, no situar.

\begin{table}[ht]
\centering
\caption{Llegadas de Bizkaibus Según los Metros de las Paradas}\label{tab:sin-metros}
\begin{tabular}{@{}lrrrrrr@{}}
\toprule
 & \multicolumn{4}{c}{Con tiempos aprendidos} & \multicolumn{2}{c}{Sin aprender} \\
\cmidrule(lr){2-5}\cmidrule(l){6-7}
Metros de las paradas & Con GPS & $\pm$1 min & $\pm$2 min & Error & $\pm$2 min & Error \\
\midrule
Del GTFS & 97,4\,\% & 55,6\,\% & 78,4\,\% & 92~s & 51,5\,\% & 192~s \\
Medidos sobre el trazado & 97,3\,\% & 55,3\,\% & 78,2\,\% & 93~s & 51,6\,\% & 192~s \\
Ninguno (recta entre paradas) & 91,9\,\% & 54,0\,\% & 76,6\,\% & 105~s & 51,2\,\% & 194~s \\
\bottomrule
\end{tabular}

\nota{Las mismas 2,08 millones de llegadas en las tres versiones. «Del GTFS»: los metros del GTFS con la corrección de
la sección anterior. «Con GPS»: las que el motor sitúa con la posición del autobús (capa~1). «Error»: error absoluto medio de la hora de llegada. «Sin aprender»: sin tiempos por tramo, factores de ritmo ni correcciones por viaje.}
\end{table}

En las cuatro redes la medida es coherente: los metros nunca bajan de una parada a la siguiente y tienen todas sus
paradas a menos de 100~m de su punto del trazado el 99,5\,\% de los viajes de Bilbobus, el 95,9\,\% de los de
Lurraldebus, el 99,6\,\% de los de Tuvisa y el 98,9\,\% de los de La Unión. Medirlos alarga la construcción de cada
red entre 0 y 9~s.

\section{Fusión de Variantes y Alias de Viaje}
\label{sec:alias}

Un mismo servicio aparece en el GTFS como varios \texttt{trip\_id} que solo se diferencian en el calendario (por
ejemplo, uno para los lunes y otro para el resto de laborables). Para no mostrar salidas duplicadas, la construcción
de la base agrupa esos \texttt{trip\_id} en un único viaje. Dos \texttt{trip\_id} se fusionan cuando coinciden en:

\begin{enumerate}
  \item la línea (\texttt{route\_id}) y el número de viaje;
  \item la secuencia de paradas (la clave del recorrido es un resumen MD5 de los identificadores de parada en orden);
  \item el grupo de calendario;
  \item la \emph{forma temporal}: para cada parada, la llegada relativa a la salida de cabecera y el tiempo de espera
        en la parada (un resumen MD5 de esa lista);
\end{enumerate}

\noindent y, además, sus salidas de cabecera son consecutivas con una separación no mayor de 90~s
(\texttt{iDepartureClusterGapSeconds}). El viaje resultante reúne los días de sus miembros (unión de las máscaras de
días de la semana y de las fechas incluidas y excluidas). Cada \texttt{trip\_id} miembro queda registrado en
\texttt{trip\_aliases} apuntando al viaje fusionado. Esa tabla es la que permite emparejar exactamente el feed en
tiempo real con el horario (sección «\nameref{sec:emparejamiento}»): el feed usa los \texttt{trip\_id} originales.

\section{Por Qué el Horario No Sirve Como Referencia Absoluta}
\label{sec:interpolado}

En el GTFS de Bizkaibus las paradas intermedias llevan \texttt{timepoint=0}. Sus horas no son compromisos del
operador, sino una interpolación lineal en la distancia entre las horas de las cabeceras. La
Figura~\ref{fig:distancia-tiempo} muestra la consecuencia. El horario es la recta $T(s)$: una velocidad constante. La
marcha real $R(s)$ es más lenta en la ciudad y mucho más rápida en la autopista. A mitad de recorrido, $R$ y $T$
pueden estar separadas muchos minutos sin que el autobús vaya «adelantado» en ningún sentido operativo.

\begin{figure}[ht]
\centering
\caption{Distancia y Tiempo de un Viaje Ilustrativo}\label{fig:distancia-tiempo}
\begin{tikzpicture}
\begin{axis}[
  width=0.88\textwidth, height=6.4cm,
  xlabel={Distancia recorrida $s$ (km)}, ylabel={Hora (min desde la salida)},
  xmin=0, xmax=30, ymin=0, ymax=62,
  axis lines=left, tick label style={font=\footnotesize}, label style={font=\small},
  legend style={font=\footnotesize,draw=none,at={(0.02,0.98)},anchor=north west}]
  \addplot[thick,color=gris,dashed,domain=0:30] {2*x};
  \addlegendentry{Horario interpolado $T(s)$}
  \addplot[very thick,color=ruta,smooth] coordinates {(0,0) (4,12) (6,17) (8,20) (14,26) (20,32) (24,40) (28,52) (30,57)};
  \addlegendentry{Marcha real $R(s)$}
  \addplot[only marks,mark=*,mark size=2.6pt,color=gps] coordinates {(14,26)};
  \draw[gps,thick,-{Stealth}] (axis cs:14,26) -- (axis cs:24,46);
  \node[font=\footnotesize,color=gps,anchor=west] at (axis cs:14.4,24.5) {GPS: $(s_p,\tgps)$};
  \draw[decorate,decoration={brace,amplitude=4pt}] (axis cs:14.3,28) -- (axis cs:14.3,26.2);
\end{axis}
\end{tikzpicture}
\nota{La recta discontinua es el horario interpolado; la curva, la marcha real. Desde el punto GPS, el método traza una paralela al horario (flecha) y la corrige tramo a tramo con los tiempos aprendidos.}
\end{figure}

\section{El Feed en Tiempo Real}
\label{sec:feed}

El feed SIRI-VM de Bizkaibus es un documento XML público que se regenera aproximadamente una vez por minuto y contiene
unos 250 vehículos en hora punta. El ritmo ha variado: en octubre de 2026 salía una versión cada 55~s, y el 5 de
noviembre, cada 64~s de mediana, con cada posición 66~s después de tomarse (sección «\nameref{sec:llegando}»). El
motor lo descarga con un tiempo máximo de 8~s y lo guarda en caché 25~s. Para cada \texttt{VehicleActivity} extrae
los campos de la Tabla~\ref{tab:feed}.

\begin{table}[ht]
\centering
\caption{Campos que el Motor Lee del Feed SIRI-VM}\label{tab:feed}
\begin{tabularx}{\textwidth}{@{}llX@{}}
\toprule
Símbolo & Elemento SIRI & Uso \\
\midrule
$\rho$ & \texttt{MonitoredVehicleJourney/VehicleJourneyRef} & Referencia del viaje; igual al \texttt{trip\_id} del GTFS. \\
$\ell,\ \nu,\ d_0^{\mathrm{feed}}$ & (extraídos de $\rho$) & Línea (sin ceros a la izquierda, como en el horario), número de viaje y salida de cabecera, por la expresión~\eqref{eq:regex}. \\
$v$ & \texttt{VehicleRef} & Identificador del vehículo. \\
$n$ & \texttt{MonitoredCall/Order} (o \texttt{VisitNumber}) & Orden de la próxima parada en el viaje. \\
$(\varphi,\lambda)$ & \texttt{VehicleLocation/Latitude}, \texttt{Longitude} & Posición GPS. \\
$\tgps$ & \texttt{LocationRecordedAtTime} (o \texttt{RecordedAtTime}) & Hora de la posición. \\
$t_{\mathrm{rec}}$ & \texttt{VehicleActivity/RecordedAtTime} & Hora del registro del vehículo. \\
$\delta_{\mathrm{decl}}$ & \texttt{Delay} (duración ISO 8601) & Retraso declarado; ausente en Bizkaibus ($0$). \\
\bottomrule
\end{tabularx}
\end{table}

\subsection{Referencia del Viaje} Solo se aceptan las actividades cuya referencia encaja con
\begin{equation}
\label{eq:regex}
\texttt{\^{}trp\_[A-Za-z]+(\textbackslash d+)\_(\textbackslash d+)\_[A-Za-z0-9]+\_(\textbackslash d+)}
\end{equation}
cuyos tres grupos son la línea $\ell$, el número de viaje $\nu$ y la salida de cabecera $d_0^{\mathrm{feed}}$ en
segundos. Por ejemplo, \texttt{trp\_A3414\_808\_OP9LSEPT\_54000\_\ldots} da $\ell=3414$, $\nu=808$ y
$d_0^{\mathrm{feed}}=54\,000$ (15:00:00). Las demás se ignoran. La línea se guarda sin ceros a la izquierda, como el
\texttt{route\_id} del horario: \texttt{trp\_A0651\_737\_\ldots} da $\ell=651$, no \texttt{0651} (el \texttt{0651} es el
relleno del nombre público \texttt{A0651}, que el horario guarda aparte en \texttt{route\_short\_name}).

\subsection{Horas} Las horas del feed son fechas ISO 8601 con zona. Se convierten a \texttt{Europe/Madrid} y a segundos
desde medianoche, $3600\,H+60\,M+S$. Si la fecha convertida no es la fecha de hoy, el valor se descarta: una posición
de ayer no dice nada del viaje de hoy. Si la hora de la posición no es válida, también se descartan la latitud y la
longitud. Una entrada sin ninguna hora de hoy, ni $\tgps$ ni $t_{\mathrm{rec}}$, no se usa: ni en la predicción, ni en
el registro de pasos, ni en el mapa (\texttt{RealtimeMatcher::hasTime}). Mientras el feed se actualiza, todas la traen: 1\,471
de 1\,471 lecturas del 6 de noviembre de 2026 entre las 23:00 y las 23:59, la última hora antes de que se detenga. De 23:58 a 05:00, en cambio, el feed de Bizkaibus no
se actualiza: repite la lectura de las 23:58 con sus horas de entonces, aunque haya autobuses nocturnos en servicio
(capturas de las noches del 31 de octubre y del 6 de noviembre de 2026). Pasada la medianoche, ninguna de esas entradas
tiene hora de hoy (sección «\nameref{sec:mas-de-24h}»).

\subsection{Retraso Declarado} El retraso, cuando existe, es una duración $\pm$\texttt{PT}$h$\texttt{H}$m$\texttt{M}$s$\texttt{S};
se convierte a $\pm(3600h+60m+s)$. En Bizkaibus no aparece y el motor de autobús no lo usa
(sección «\nameref{sec:otras-redes}»).

Las entradas se agrupan por la clave «$\ell|\nu$», porque puede haber varios vehículos con la misma línea y número de
viaje en distintos días de calendario.

\subsection{Lo que el Feed No Dice} El feed solo trae posiciones: no publica cancelaciones, ni qué vehículo hará cada
viaje, ni refuerzos fuera del horario. Con las capturas del 3 de noviembre de 2026 se midieron cuatro consecuencias:

\begin{itemize}
  \item \emph{Viajes sin autobús.} En cuatro momentos del día, 947 viajes debían estar en marcha según el horario y
        187 (el 20\,\%) no tenían autobús en el feed. Siguiendo esos viajes en las capturas, la mayoría no falta: el
        12\,\% del total ya había llegado al final antes de su hora (el horario deja margen al final del recorrido),
        un 2\,\% aparece después con unos 15~min de retraso (mediana), otro 2\,\% desaparece un rato y vuelve, y otro
        2\,\% desaparece a mitad y no vuelve. Solo el 2\,\% no aparece en todo el día; desde fuera no se distingue
        un autobús sin GPS de un viaje que no se hizo.
  \item \emph{El número de viaje no es el vehículo.} De los viajes seguidos de una misma clave «$\ell|\nu$», solo
        el 11\,\% los hizo el mismo autobús, y el GTFS no trae \texttt{block\_id}. No se puede saber de antemano qué
        autobús hará el próximo viaje, ni arrastrar a él el retraso del anterior (sección «\nameref{sec:descartadas}»).
  \item \emph{Orden y parada.} El orden $n$ y la parada \texttt{StopPointRef} coinciden en el 97{,}9\,\% de los
        casos. En unas pocas líneas (3129, 2610, 2611 y 652) el operador cuenta una o varias paradas más que el GTFS;
        el motor usa $n$, y fiarse de la parada no mejoró la medida.
  \item \emph{Dónde está el autobús.} El feed no dice si un autobús está en una parada. \texttt{StopPointRef} es
        siempre la próxima: mientras espera en la primera parada de su viaje, apunta a la segunda (359 de 360 casos en
        nueve capturas del 3 de noviembre). \texttt{ArrivalStatus} vale \texttt{arrived} en los 2\,142 autobuses
        mirados, parados o en marcha, y no hay \texttt{PreviousCalls}. Solo el GPS lo dice (sección
        «\nameref{sec:cabecera}»).
\end{itemize}

\section{Notación}
\label{sec:notacion}

\begin{definicion}[Viaje]
Un viaje $J$ es una secuencia de $N\ge2$ paradas $p_1,\dots,p_N$. Para cada $i$ se conocen la llegada y la salida
programadas $a_i\le d_i$, las coordenadas $(\varphi_i,\lambda_i)$, el identificador $u_i$ de la parada y,
opcionalmente, la distancia $s_i$ a lo largo del trazado, con $s_1\le s_2\le\dots\le s_N$. El orden de la parada
$p_i$ es $i$ (\texttt{seq\_order}). La salida de cabecera es $d_1$. El trazado, si existe, es una polilínea
$\Sigma=\big((\varphi_j,\lambda_j,\sigma_j)\big)_{j=1}^{M}$ con $\sigma_j$ no decreciente.
\end{definicion}

\begin{definicion}[Tramo]
El tramo $i$ de $J$ es el par de paradas consecutivas $(p_i,p_{i+1})$. Su identificador es «$u_i>u_{i+1}$», su
duración programada es $a_{i+1}-a_i$ y su longitud, si se conocen las distancias, es $s_{i+1}-s_i$. Dos viajes de
líneas distintas que pasan por las mismas dos paradas comparten el tramo.
\end{definicion}

\begin{definicion}[Entrada del feed]
Una entrada es la tupla $e=(\rho,\allowbreak\ell,\allowbreak\nu,\allowbreak d_0^{\mathrm{feed}},\allowbreak v,\allowbreak n,\allowbreak\varphi,\allowbreak\lambda,\allowbreak\tgps,\allowbreak t_{\mathrm{rec}})$ de la
Tabla~\ref{tab:feed}, donde cualquiera de $n,\varphi,\lambda,\tgps,t_{\mathrm{rec}}$ puede faltar.
\end{definicion}

\begin{definicion}[Franja]
\label{def:franja}
La franja de una fecha $D$ y un instante $t$ es el par
\begin{gather*}
\franja(D,t)=\big(\text{tipo}(D),\ h(t)\big),\\
\text{tipo}(D)=\begin{cases}\texttt{fin\_semana}&\text{si $D$ es sábado o domingo},\\ \texttt{laborable}&\text{en otro caso,}\end{cases}
\end{gather*}
con $h=\lfloor t/3600\rfloor \bmod 24$ y
\[
h(t)=\begin{cases}
\texttt{madrugada} & 0\le h<6\\
\texttt{punta\_manana} & 6\le h<9\\
\texttt{manana} & 9\le h<13\\
\texttt{mediodia} & 13\le h<16\\
\texttt{punta\_tarde} & 16\le h<20\\
\texttt{noche} & 20\le h<24.
\end{cases}
\]
Se escribe como texto «\texttt{tipo.nombre}», por ejemplo \texttt{laborable.punta\_manana}. Hay, por tanto, doce
franjas.
\end{definicion}

\begin{observacion}
La fecha que se usa es la del día de servicio, no la del reloj. Un instante de 25:00 ($90\,000$~s) de un miércoles
es la franja \texttt{laborable.madrugada} del miércoles, aunque en el calendario ya sea jueves. En los tiempos por
tramo, los festivos entre semana se tratan como laborables: se comprobó que su tráfico se parece más al de un laborable
que al de un fin de semana (sección «\nameref{sec:descartadas}»). La capa 3 es la excepción: para ella el festivo
cuenta como fin de semana en todo (franja de los tramos, factor de ritmo, parte aprendida y corrección por viaje),
porque así se midió mejor (sección «\nameref{sec:capa3}»).
\end{observacion}

\begin{definicion}[Hora actual]
$\tnow$ es la hora del servidor en el momento de responder, en segundos desde la medianoche de Madrid del día actual.
Para un viaje del día anterior se usa $\tnow+86\,400$, en las horas de su horario (sección «\nameref{sec:mas-de-24h}»).
\end{definicion}
